\section{Core Ablations and Screening Studies}
\label{app:design}

\paragraph{Final-scale confirmation.}
Table~\mainref{tab:design} contains four independently trained controls of the
24-layer final model, with three training seeds each. They keep the final data,
initialization recipe, training exposure, and other components fixed
(Appendix~\ref{app:recipe}). Evaluation uses 1{,}400 held-out scenes under the
real-time protocol at 0 and 150\,ms added delay, with shared scene keys and
inference seeds. The full-model reference is reused from the main evaluation.
The four controls are:
\begin{itemize}[leftmargin=*,topsep=2pt,itemsep=1pt]
\item \textbf{No target history:} retain the current image, current target pose,
and eight robot states; remove past target poses and their velocity and
acceleration features.
\item \textbf{No future supervision:} disable SE(3), scene-flow and visibility,
future depth, DINO, and dynamic-RGB losses; retain current-target RGB
reconstruction and the action loss.
\item \textbf{No scene-flow supervision:} disable only scene-flow and visibility
losses; retain dedicated SE(3) and all appearance targets.
\item \textbf{No dedicated SE(3):} disable only the target-pose forecast loss;
retain all appearance targets and the same $14^2$ scene-flow queries over
the target, robot, and surroundings.
\end{itemize}
The last control assigns the target the same coarse flow representation as the
rest of the scene; it introduces no new target-track head. Since scene flow
includes the target and robot, removing it does not isolate environment
information. Results are pending: \fillin{complete the four controls and their
paired intervals before drawing final-scale conclusions}.

\paragraph{Screening scope.}
The broader studies used 4 layers of width 768 and 4 epochs, generally 40{,}104
updates at global batch size 16. HISTORY1 used a frozen DINO-L encoder and one
training seed, with a first-epoch cosine schedule followed by a low learning
rate. FUTURE0 used the same depth and width and one seed, but a cosine schedule
over all four epochs. Q3-INFO0 used MAE and three seeds for each of R0--R5;
BACKBONE0 compared its two backbones with three seeds. The added AUX0 appearance
conditions used one seed. These groups are not one matched final-model
comparison. Development results motivate design choices; they retain their
original split and recipe when cited. The retained screening tables describe
completed conditions at their original scale; they do not add 24-layer training
arms to the core confirmation.

\subsection{Shared Components}
\label{app:design_shared}
The encoder and backbone comparisons belong to the screening scale. They
compare a trained CNN, frozen DINOv3 ViT-S/16 and ViT-L/16 \cite{dinov3}, and
MAE ViT-B/16 \cite{mae}, with their original group-specific recipes and seed
counts. The MAE reference and the BACKBONE0 comparison have three seeds; this
does not make every encoder or auxiliary-target condition a three-seed study.
The final model uses MAE; DINO-L remains a future feature target. Table~\ref{tab:supp_shared}
retains the screening comparison layout.
% Transcribed from Research docs/experiments/SCALE0/SCALE0.md and AUX0/results/comparison.md (8-GPU encoder runs), BACKBONE0/BACKBONE0.md (matched parameter sharing). SR mean/sample SD from three seed counts /1400.
\begin{table}[!htb]
\centering\footnotesize
\caption{\tablelang{\textbf{Encoder and parameter sharing.} Three seeds, 1{,}400 scenes per seed; direct 3D, SE(3), and all appearance targets. SR: seed mean $\pm$ sample s.d.\ (\%). Panel (b) fixes MAE-B, reusing (a)'s shared-attention runs. Params: trainable backbone; $\star$: selected. Blue/bold and light blue mark the best and second point estimates within each panel and outcome column, not statistical significance.}{\textbf{编码器与参数共享。}三个种子，各评测 1{,}400 场；直接三维、SE(3) 与全部外观目标。SR 为种子均值 $\pm$ 样本标准差（\%）。面板 (b) 固定 MAE-B，复用 (a) 的共享注意力结果。参数量仅计可训练主干；$\star$ 表示选定配置。 各面板和结果列中，粗体加蓝底与浅蓝底分别标出最佳、次佳点估计，不表示统计显著性。}}
\label{tab:supp_shared}
\begin{tabularx}{\columnwidth}{@{}Xlc@{}}
\toprule
\tablelang{Configuration}{配置} & \tablelang{Status / params}{状态／参数量} & SR (\%) \\
\midrule
\multicolumn{3}{@{}l}{\tablelang{\textit{(a) Observation encoder}}{\textit{(a) 观测编码器}}}\\
CNN & \tablelang{trainable}{可训练} & \secondcell{$62.86\pm2.64$} \\
DINOv3 ViT-S/16 & \tablelang{frozen}{冻结} & $61.69\pm2.00$ \\
DINOv3 ViT-L/16 & \tablelang{frozen}{冻结} & $62.21\pm1.82$ \\
MAE ViT-B/16\sel & \tablelang{frozen}{冻结} & \bestcell{$66.76\pm3.84$} \\
\midrule
\multicolumn{3}{@{}l}{\tablelang{\textit{(b) World/action parameter sharing}}{\textit{(b) 世界／动作参数共享}}}\\
\tablelang{Shared attention\sel}{共享注意力\sel} & 47.266M & \bestcell{$66.76\pm3.84$} \\
\tablelang{Separate parameters}{独立参数} & 47.272M & \secondcell{$62.86\pm2.98$} \\
\bottomrule
\end{tabularx}
\end{table}


\subsection{History}
\label{app:design_remember}
HISTORY1 screened current-state input, target motion history, RGB windows,
Farneb\"ack flow \cite{farneback}, flow with motion statistics
\cite{dynamicwam}, codec motion vectors, video-VAE latents, and compressed
windows. Its nine trained conditions use one seed each. Table~\ref{tab:supp_remember}
reports these nine completed conditions. Only the no-target-history control is repeated at the final scale
in the core confirmation.

% Transcribed from Research HISTORY1/HISTORY1.md, completed nine-condition seed-0 study; original 4-layer synchronous evaluation.
\begin{table}[t]
\centering\footnotesize
\caption{\tablelang{\textbf{History screening.} Four layers, width 768, frozen DINO-L, one training seed, and the original synchronous evaluation on 1,400 scenes. The first eight conditions receive the current target pose; the final condition adds target-motion history. These are not final-scale real-time results.}{\textbf{历史筛选。}4层、宽768、冻结DINO-L、一个训练种子，按原同步协议评测1,400个场景。前八项接收当前目标位姿，最后一项另加目标运动历史。这些不是最终规模的实时结果。}}
\label{tab:supp_remember}
\begin{tabular}{@{}lrr@{}}
\toprule
\tablelang{History input}{历史输入} & \tablelang{Success / 1,400}{成功数 / 1,400} & SR (\%) \\
\midrule
\tablelang{Current image}{当前图像} & 748 & 53.43 \\
\tablelang{RGB window}{RGB窗口} & 662 & 47.29 \\
\tablelang{Optical flow}{光流} & 720 & 51.43 \\
\tablelang{Flow + motion statistics}{光流＋运动统计} & 723 & 51.64 \\
\tablelang{Codec motion vectors}{编码器运动向量} & 733 & 52.36 \\
\tablelang{Video-VAE latent}{视频VAE隐变量} & 710 & 50.71 \\
\tablelang{Compressed window (age)}{压缩窗口（按时间）} & 634 & 45.29 \\
\tablelang{Compressed window (uniform)}{压缩窗口（均匀）} & 670 & 47.86 \\
\midrule
\tablelang{Current image + target-motion history}{当前图像＋目标运动历史} & 824 & 58.86 \\
\bottomrule
\end{tabular}
\end{table}


\paragraph{Stationary-history intervention.}
The completed intervention in Table~\ref{tab:supp_interventions} uses the frozen
4-layer HISTORY1 checkpoint, one training seed, and the original synchronous
executor. It replaces past target poses by the current pose before recomputing
kinematics; current pose, image, robot state, timestamps, and validity are retained.
The original history result is reused. This intervention changes the input
distribution and does not replace the separately trained final-scale no-history
control.
% Transcribed from Research HISTORY1/HISTORY1.md, frozen GT-E intervention (2026-10-18), original 4-layer seed-0 synchronous evaluation.
\begin{table}[t]
\centering\footnotesize
\caption{\tablelang{\textbf{Stationary-history intervention.} One frozen 4-layer HISTORY1 checkpoint, training seed 0, and its original synchronous protocol. Past target poses are replaced by the current pose and kinematics recomputed; the unchanged-history result is reused. These 1,400 original scenes are separate from the final real-time confirmation.}{\textbf{静止历史干预。}使用同一个冻结的4层 HISTORY1 检查点、训练种子0及其原同步协议。将过去目标位姿替换为当前位姿，再重算运动学；真实历史结果直接复用。这1,400个原场景与最终实时确认实验分开报告。}}
\label{tab:supp_interventions}
\begin{tabular}{@{}lrr@{}}
\toprule
\tablelang{Target history}{目标历史} & \tablelang{Success / 1,400}{成功数 / 1,400} & SR (\%) \\
\midrule
\tablelang{Unchanged}{真实历史} & 824 & 58.86 \\
\tablelang{Stationary replacement}{静止历史替换} & 9 & 0.64 \\
\bottomrule
\end{tabular}
\end{table}


\paragraph{Pose noise.}
The separate robustness evaluation freezes the model and reuses the real-time
clean reference. It adds three conditions: nominal synthetic temporally
correlated pose noise, twice its amplitude, and nominal noise with contiguous
missing observations. The perturbation applies to current and historical
estimates before history caching and kinematics. Actual tracker calibration
and closed-loop tracking remain pending (Table~\ref{tab:supp_tracker}).
No history-only noise crossover or additional visual-history training is required
for this three-condition evaluation.
The nominal process combines a fixed per-episode bias and a stationary correlated
component, each with per-axis standard deviations of 5\,mm and $1^\circ$ for
translation and rotation. The correlation time is 200\,ms at 50\,ms sampling;
the doubled condition scales both components. Dropout lasts four ticks
(200\,ms), with geometrically distributed valid gaps of mean 76 ticks, giving
a long-run missing fraction of 5\%. These parameters are set before evaluation.
All three final-model training seeds use the same matched scenes at zero added
latency; three noise conditions are not three training seeds.

\subsection{Forecast}
\label{app:design_predict}
FUTURE0 screened fifteen target configurations with one seed each.
Table~\ref{tab:target_shapes} defines their representations, and
Table~\ref{tab:supp_predict} reports the completed conditions. P09 uses target
2D tracks; it is distinct from the final-scale uniform 3D-flow control.
The table also reports AUX0's four MAE appearance-target conditions and their
existing all-target reference, all at seed 1. These use a different encoder and
recipe from FUTURE0; their scores are compared within each group. They cannot be
pooled with Q3's three-seed coupling results or presented as final-scale evidence
that geometry outperforms appearance.

\begin{table*}[t]
\centering\small
\caption{\tablelang{Future targets in the representation screen. Shape excludes batch/query axes.
Displacements are cumulative from the current state except codec vectors,
which retain the decoder's native correspondence. All targets come from
the recorded demonstration.}{表示筛选中的未来监督目标。张量形状不含 batch 和 query 轴。除保留解码器原生对应关系的 codec 向量外，位移均从当前状态累计。所有监督目标均来自记录的演示。}}
\label{tab:target_shapes}
\begin{tabularx}{\textwidth}{@{}llX@{}}
\toprule
\tablelang{Target}{目标} & \tablelang{Shape}{形状} & \tablelang{Support, units, and validity}{几何支持、单位与有效条件} \\
\midrule
\tablelang{Scene flow}{场景流} & $196\times8\times3$ & \tablelang{Fixed $14\times14$ patch centers; base-frame meters; valid source depth/geometry and finite future displacement.}{固定 $14\times14$ patch 中心；以基座坐标系中的米为单位；源深度／几何有效，未来位移有限。} \\
\tablelang{Target pose}{目标位姿} & $8\times6$ & \tablelang{Cumulative translation (m) and left-relative rotvec (rad); valid object pose and query.}{累计平移 (m) 和左相对旋转向量 (rad)；物体位姿和 query 有效。} \\
P04: Wan & $48\times2\times12\times10$ & \tablelang{Frozen Wan2.2 VAE posterior mean and native normalization; future groups end at 200/400\,ms; valid query.}{冻结的 Wan2.2 VAE 后验均值及原生归一化；未来分组终点为 200/400\,ms；query 有效。} \\
\tablelang{P05: dense 2-D}{P05：稠密二维} & $8\times2\times224\times224$ & \tablelang{Fixed source pixels; displacement divided by 224; valid source and finite future projection.}{固定源像素；位移除以 224；源像素有效，未来投影有限。} \\
\tablelang{P06: codec}{P06：codec} & $8\times2\times64\times64$ & \tablelang{Each future/current pair is encoded as P/I frames; native vectors divided by 64; actual valid vector regions.}{每对未来／当前帧编码为 P/I 帧；原生向量除以 64；使用实际有效向量区域。} \\
\tablelang{P07: scene tracks}{P07：场景轨迹} & $196\times8\times3$ & \tablelang{Same patch source points; two displacement channels divided by 224 and one visibility logit; valid correspondences.}{相同的 patch 源点；两个位移通道除以 224，另加一个可见性 logit；对应关系有效。} \\
\tablelang{P08: object DINO}{P08：物体 DINO} & $8\times1024$ & \tablelang{Object-mask-area-weighted patch pooling; channel layer normalization; positive mask area.}{按物体掩码面积加权汇聚 patch；沿通道作层归一化；掩码面积为正。} \\
\tablelang{P09: object tracks}{P09：物体轨迹} & $32\times8\times3$ & \tablelang{Deterministic farthest-point sampling in the current target mask; displacement and visibility as P07; padded points invalid.}{在当前目标掩码内确定性地作最远点采样；位移和可见性同 P07；填充点无效。} \\
\tablelang{$A_\alpha$: magnitude}{$A_\alpha$：幅值} & $196\times8$ & \tablelang{Cumulative flow norm (m); same correspondence mask as flow.}{累计流的范数 (m)；对应关系掩码与场景流相同。} \\
\tablelang{P12: depth}{P12：深度} & $8\times1\times224\times224$ & \tablelang{Metric future depth; finite positive pixels.}{未来深度的米制值；像素值为有限正数。} \\
P13: DINO & $8\times196\times1024$ & \tablelang{Frozen DINOv3-L/16 patches, channel layer normalization; valid query/horizon/patch.}{冻结的 DINOv3-L/16 patch，沿通道作层归一化；query、时域和 patch 均有效。} \\
\tablelang{P14: dynamic RGB}{P14：动态 RGB} & $8\times3\times224\times224$ & \tablelang{RGB divided by 255; visible surfaces moving more than 1\,mm in world coordinates, including robot and other geometry.}{RGB 除以 255；世界坐标中移动超过 1\,mm 的可见表面，包括机器人和其他几何体。} \\
\bottomrule
\end{tabularx}
\end{table*}

% Transcribed from tables/future_complete.tex (Research FUTURE0 accepted summaries) and Research AUX0/AUX0.md MAE-B seed-1 results (2026-10-22).
\begin{table*}[p]
\centering\small
\caption{\tablelang{\textbf{Forecast targets} at the screening scale. The first group uses a frozen DINO-L encoder and seed 0, with its 140 development scenes and the 1{,}260 remaining scenes reported separately. The second group uses the MAE encoder and seed 1 on 1{,}400 scenes; it keeps scene flow, SE(3), and the current-target RGB reconstruction and removes future appearance targets. Compare conditions within each group.}{\textbf{预测目标}（筛选规模）。第一组使用冻结 DINO-L 编码器与种子 0，其 140 个开发场景与其余 1{,}260 个场景分别报告。第二组使用 MAE 编码器与种子 1，在 1{,}400 个场景上评测；它保留场景流、SE(3) 与当前目标 RGB 重建，并移除未来外观目标。请在组内比较各条件。}}
\label{tab:supp_predict}
\setlength{\tabcolsep}{4pt}
\begin{tabular*}{\textwidth}{@{\extracolsep{\fill}}llrrr@{}}
\toprule
\tablelang{ID}{编号} & \tablelang{Future supervision}{未来监督} & \tablelang{Dev.\ success / 140}{开发集成功数 / 140} & \tablelang{Success / 1,260}{成功数 / 1,260} & SR$_{1260}$ (\%) \\
\midrule
P00 & \tablelang{None}{无} & 88 & 791 & 62.78 \\
P01 & \tablelang{3D scene flow}{三维场景流} & 95 & 782 & 62.06 \\
P02 & \tablelang{SE(3)}{SE(3)} & 92 & 789 & 62.62 \\
P03 & \tablelang{Scene flow + SE(3)}{场景流＋SE(3)} & 93 & 786 & 62.38 \\
P04 & \tablelang{Video-VAE latent + SE(3)}{视频VAE隐变量＋SE(3)} & 92 & 783 & 62.14 \\
P05 & \tablelang{Dense 2D displacement + SE(3)}{稠密二维位移＋SE(3)} & 87 & 746 & 59.21 \\
P06 & \tablelang{Codec motion + SE(3)}{编码器运动向量＋SE(3)} & 86 & 795 & 63.10 \\
P07 & \tablelang{Sparse scene 2D tracks + SE(3)}{稀疏场景二维轨迹＋SE(3)} & 90 & 807 & 64.05 \\
P08 & \tablelang{Scene flow + object DINO}{场景流＋物体DINO} & 89 & 794 & 63.02 \\
P09 & \tablelang{Scene flow + target 2D tracks}{场景流＋目标二维轨迹} & 88 & 808 & 64.13 \\
$A_\alpha$ & \tablelang{P03 + motion magnitude}{P03＋运动幅值} & 89 & 778 & 61.75 \\
P12 & \tablelang{P03 + depth}{P03＋深度} & 90 & 791 & 62.78 \\
P13 & \tablelang{P03 + DINO}{P03＋DINO} & 92 & 788 & 62.54 \\
P14 & \tablelang{P03 + dynamic RGB}{P03＋动态RGB} & 89 & 790 & 62.70 \\
P15 & \tablelang{P03 + depth / DINO / dynamic RGB}{P03＋深度／DINO／动态RGB} & 94 & 802 & 63.65 \\
\midrule
\multicolumn{2}{@{}l}{\tablelang{AUX0 MAE-B: future appearance auxiliaries}{AUX0 MAE-B：未来外观辅助}} & \multicolumn{2}{c}{\tablelang{Success / 1,400}{成功数 / 1,400}} & SR$_{1400}$ (\%) \\
\midrule
\multicolumn{2}{@{}l}{\tablelang{None}{无}} & \multicolumn{2}{c}{894} & 63.86 \\
\multicolumn{2}{@{}l}{\tablelang{Depth only}{仅深度}} & \multicolumn{2}{c}{867} & 61.93 \\
\multicolumn{2}{@{}l}{\tablelang{DINO only}{仅DINO}} & \multicolumn{2}{c}{946} & 67.57 \\
\multicolumn{2}{@{}l}{\tablelang{Dynamic RGB only}{仅动态RGB}} & \multicolumn{2}{c}{926} & 66.14 \\
\multicolumn{2}{@{}l}{\tablelang{All three (existing reference)}{三者全部（已有参照）}} & \multicolumn{2}{c}{995} & 71.07 \\
\bottomrule
\end{tabular*}
\end{table*}


\paragraph{Role coverage.}
The core no-SE(3) control retains the original $14^2$ scene flow, and the
no-flow control retains SE(3) and all appearance targets. Since the flow covers
target, robot, and context, the latter tests scene-flow supervision as a whole.
The study makes no claim that a coarse grid matches denser grids.

\paragraph{Extrapolation and forecast errors.}
Constant-velocity and constant-acceleration estimates from causal pose history
can be compared offline with learned forecasts at the logged execution times.
Table~\ref{tab:supp_extrapolation} compares their position errors on the same
held-out observations at a fixed 400\,ms horizon. CV and CA use only causal pose
history; the learned head reads the frozen full-model states. These are offline
prediction diagnostics and introduce no controller or training arm. They do not
show that the deployed policy reads decoded forecasts. The separately trained
CV-input policy remains in Table~\mainref{tab:main}; no additional CV/CA-input
training combinations are added.
% Planned results: offline causal CV/CA and frozen learned-pose-head errors; no new controller or training arm.
\begin{table}[!htb]
\centering\footnotesize
\caption{\tablelang{\textbf{Offline prediction errors.} Target-position error at 400\,ms (mm) on the same held-out observations. CV and CA are analytic extrapolations from causal pose history; the learned head reads frozen full-model states.}{\textbf{离线预测误差。}在同一批留出观测上计算400\,ms处的目标位置误差（mm）。CV和CA是由因果位姿历史计算的解析外推；学习位姿头读取冻结完整模型的状态。}}
\label{tab:supp_extrapolation}
\setlength{\tabcolsep}{2.6pt}
\begin{tabular}{@{}lcccc@{}}
\toprule
 & \tablelang{Const.\ vel.}{匀速} & \tablelang{Harmonic}{谐波} & \tablelang{Roll/slide}{滚动/滑动} & \tablelang{Impact}{撞击} \\
\midrule
\tablelang{CV extrapolation}{匀速外推} & \resultplaceholder & \resultplaceholder & \resultplaceholder & \resultplaceholder \\
\tablelang{CA extrapolation}{匀加速外推} & \resultplaceholder & \resultplaceholder & \resultplaceholder & \resultplaceholder \\
\tablelang{Learned pose head}{学习的位姿头} & \resultplaceholder & \resultplaceholder & \resultplaceholder & \resultplaceholder \\
\bottomrule
\end{tabular}
\end{table}


\subsection{Use of the Forecast and Backbone}
\label{app:design_act}
BACKBONE0 compared shared attention with stream-specific feed-forward networks
against separate world/action parameter groups at about 47.3M backbone
parameters. Its three-seed results belong to this 4-layer setting.

Q3-INFO0 trained R0--R5 at the same screening scale with three seeds each
(Appendix~\ref{app:coupling}). Its R1 contains world queries with no edge to the
actions, whereas the final model reads out from summary tokens without these
queries. The Q3 comparisons therefore do not directly measure the final
model. Coupling effects changed sign across seeds; this supports the choice
of a simpler deployment but does not establish statistical equivalence.
R1 minus R0 was \fillin{3.14} percentage points on average, with seed-wise
differences of \fillin{$-$1.29}, \fillin{+9.64}, and \fillin{+1.07}.
The core confirmation does not repeat R5 or the other coupling designs at
24 layers.

% Planned results: Q3-INFO0 4x768, MAE, 4 epochs, three training seeds per R0-R5. Screening layout; no final-scale R5 run or equivalence claim.
\begin{table}[t]
\centering\footnotesize
\caption{\tablelang{\textbf{Couplings R0--R5 and backbones} at the screening scale (MAE encoder, three training seeds each; Appendix~\ref{app:design_act}).}{\textbf{耦合方式 R0--R5 与主干}（筛选规模：MAE 编码器，各三个训练种子；附录~\ref{app:design_act}）。}}
\label{tab:supp_act}
\setlength{\tabcolsep}{4.4pt}
\begin{tabular}{@{}lccc@{}}
\toprule
\tablelang{Coupling}{耦合} & \tablelang{SR mean}{SR 均值} & \tablelang{SR std.}{SR 标准差} & p95 (ms) \\
\midrule
R0 \tablelang{(future losses off)}{（关闭未来损失）} & \resultplaceholder & \resultplaceholder & \resultplaceholder \\
R1 \tablelang{(queries, no edge)}{（查询，无连接）} & \resultplaceholder & \resultplaceholder & \resultplaceholder \\
R2 \tablelang{(world$\to$action)}{（世界$\to$动作）} & \resultplaceholder & \resultplaceholder & \resultplaceholder \\
R3 \tablelang{(action$\to$world)}{（动作$\to$世界）} & \resultplaceholder & \resultplaceholder & \resultplaceholder \\
R4 \tablelang{(bidirectional)}{（双向）} & \resultplaceholder & \resultplaceholder & \resultplaceholder \\
R5 \tablelang{(decoded geometry)}{（解码几何）} & \resultplaceholder & \resultplaceholder & \resultplaceholder \\
\bottomrule
\end{tabular}
\end{table}


\paragraph{Deployment check.}
On the final checkpoint, running unused forecast heads can check equality of
the action output for fixed inputs and flow-matching noise and quantify the
added computation. Since the real-time executor advances physics during this
computation, equal outputs at fixed inputs do not imply equal closed-loop
success. No forecast-swap, time-shift, or oracle-input matrix is added for R2/R5.
